%=============================================================================
% SPECTRAL ACTION COMPUTATION IN DFD: THE COMPLETE CHAIN TO kappa_{U(1)} = 7.26999
% Cross-reference document extracting the full spectral action derivation
% from the DFD Unified Theory (v108) and Ab Initio paper
% Compiled 2026-03-22
%=============================================================================
\documentclass[12pt,a4paper]{article}
\usepackage{amsmath,amssymb,amsthm}
\usepackage{booktabs}
\usepackage{geometry}
\geometry{margin=2.5cm}
\usepackage{hyperref}
\usepackage{tcolorbox}

\newtheorem{theorem}{Theorem}
\newtheorem{lemma}{Lemma}
\newtheorem{corollary}{Corollary}
\newtheorem{definition}{Definition}
\newtheorem{remark}{Remark}

\title{The Spectral Action Computation in DFD:\\[6pt]
\large How $\kappa_{U(1)} = 7.26999\ldots$ Produces
$\alpha^{-1} = 137.03599985$}
\author{Cross-reference extracted from G.\ Alcock,\\
\textit{Density Field Dynamics: A Complete Unified Theory} (v108)\\
and \textit{Ab Initio Derivation of the Fine Structure Constant} (v12)}
\date{March 2026}

\begin{document}
\maketitle

\begin{abstract}
This document collates the complete spectral action computation chain
in Density Field Dynamics (DFD) that produces the U(1) stiffness
coefficient $\kappa_{U(1)} = 7.26999\ldots$ and thence the fine-structure
constant $\alpha^{-1} = 137.03599985$ (residual $-0.006$~ppm from experiment).
The derivation has two layers: (1)~an exact electroweak mixing formula
$\alpha^{-1} = 6\pi\,\kappa_{U(1)}$, and (2)~a spectral action / heat-kernel
computation of $\kappa_{U(1)}$ from the $a_4$ Seeley--DeWitt coefficient on
the Toeplitz-truncated internal space $\mathbb{C}P^2 \times S^3$ at
truncation level $d = 64$.
All source locations in the DFD codebase are cited.
\end{abstract}

\tableofcontents
\newpage

%=============================================================================
\section{Overview: The Two-Layer Structure}
\label{sec:overview}
%=============================================================================

The DFD derivation of the fine-structure constant proceeds in two
independent layers:

\begin{tcolorbox}[colback=blue!5!white, colframe=blue!50!black,
  title=Two-Layer Architecture]
\textbf{Layer 1 (Exact):} Electroweak mixing gives
\[
\alpha^{-1} = 6\pi\,\kappa_{U(1)}\,.
\]
This is derived from $1/e^2 = 1/g_1^2 + 1/g_2^2$ with the DFD
stiffness ratio $\kappa_{U(1)}/\kappa_{SU(2)} = 1/2$ (Frame Stiffness
Theorem).

\medskip
\textbf{Layer 2 (Numerical):} The spectral action on
$\mathbb{C}P^2 \times S^3$ with Toeplitz truncation at $d = 64$
predicts
\[
\kappa_{U(1)} = 7.26999\ldots
\]
from the $a_4$ Seeley--DeWitt coefficient, using eight locked
geometric/SM inputs (Table~XXVIII of the Unified Theory).
\end{tcolorbox}

\noindent
The two layers combine to give:
\begin{equation}
\boxed{\alpha^{-1} = 6\pi \times 7.26999\ldots = 137.03599985}
\end{equation}
with residual $-0.006$~ppm from the experimental value
$\alpha^{-1}_{\mathrm{exp}} = 137.035999084(21)$.

\paragraph{Source files.}
\begin{itemize}
\item Layer 1: \texttt{section\_alpha\_locked.tex} (Sec.~8C),
  \texttt{appendix\_F.tex} (gauge emergence),
  \texttt{appendix\_G.tex} ($\alpha$-relation derivations)
\item Layer 2: \texttt{section\_alpha\_locked.tex} (Sec.~8C),
  \texttt{appendix\_K.tex} (microsector physics),
  \texttt{Alpha\_Inverse\_Formula.tex} (explicit chain)
\item Lattice verification: \texttt{appendix\_K.tex} (Sec.~K.2)
\end{itemize}

%=============================================================================
\section{Layer 1: The Exact Electroweak Formula $\alpha^{-1} = 6\pi\,\kappa_{U(1)}$}
\label{sec:layer1}
%=============================================================================

%---------------------------------------------------------------------
\subsection{Gauge Emergence and Frame Stiffness}
%---------------------------------------------------------------------

In DFD, gauge fields emerge as Berry connections on the internal
manifold $\mathcal{M}_7 = \mathbb{C}P^2 \times S^3$. The gauge group
$SU(3) \times SU(2) \times U(1)$ arises from the unique minimal
block partition $(3,2,1)$ of $\mathbb{C}^6$ (Proposition~F.1 in
\texttt{appendix\_F.tex}).

The gauge-kinetic Lagrangian for sector $r$ is:
\begin{equation}
\mathcal{L}^{(r)}_{\mathrm{stiff}}
  = -\frac{\kappa_r}{2}\,\mathrm{Tr}\,F^{(r)}_{ij} F^{(r)\,ij}\,,
\end{equation}
where $\kappa_r$ is the frame stiffness coefficient.

\begin{theorem}[Frame Stiffness from Ricci Curvature; Thm.~F.6]
\label{thm:frame-stiffness}
The frame stiffness satisfies $\kappa_r = n_r \cdot \kappa_0$,
where $n_r$ is the block dimension in the $(3,2,1)$ partition.
\end{theorem}

\noindent
This gives the stiffness ratio:
\begin{equation}
\frac{\kappa_{U(1)}}{\kappa_{SU(2)}} = \frac{n_1}{n_2} = \frac{1}{2}\,.
\end{equation}

%---------------------------------------------------------------------
\subsection{Wilson Normalization and Electroweak Mixing}
%---------------------------------------------------------------------

The gauge couplings are extracted via Wilson normalization:
\begin{equation}
g_1^2 = \frac{1}{\kappa_{U(1)}}\,,\qquad
g_2^2 = \frac{4}{\kappa_{SU(2)}}\,.
\end{equation}

The electromagnetic coupling satisfies:
\begin{equation}
\frac{1}{e^2} = \frac{1}{g_1^2} + \frac{1}{g_2^2}
= \kappa_{U(1)} + \frac{\kappa_{SU(2)}}{4}\,.
\end{equation}

Substituting $\kappa_{SU(2)} = 2\kappa_{U(1)}$:
\begin{equation}
\frac{1}{e^2} = \kappa_{U(1)} + \frac{2\kappa_{U(1)}}{4}
= \frac{3}{2}\,\kappa_{U(1)}\,.
\end{equation}

Since $\alpha = e^2/(4\pi)$:
\begin{equation}
\boxed{\alpha^{-1} = \frac{4\pi}{e^2}
= 4\pi \cdot \frac{3}{2}\,\kappa_{U(1)}
= 6\pi\,\kappa_{U(1)}\,.}
\end{equation}

This is exact within the DFD gauge emergence framework.
The required value is:
\begin{equation}
\kappa_{U(1)} = \frac{\alpha^{-1}_{\mathrm{exp}}}{6\pi}
= \frac{137.035999084}{6\pi} = 7.269986\ldots
\end{equation}

%---------------------------------------------------------------------
\subsection{Connection to the Weinberg Angle}
%---------------------------------------------------------------------

The stiffness ratio also determines the Weinberg angle. From the
partition weights $(3,2,1)$ with canonical hypercharge normalization
$\mathrm{Tr}(Y^2) = 10/3$ per generation:
\begin{equation}
\sin^2\theta_W = \frac{g'^2}{g'^2 + g^2}
= \frac{3}{13} = 0.2308\,,
\end{equation}
matching the measured value $0.23122 \pm 0.00004$ to within $0.2\%$
(the offset is consistent with radiative corrections from tree level
to $\overline{\mathrm{MS}}$).

Source: \texttt{supplementary/PRL\_SM\_from\_Topology.tex}, Sec.~2.

%=============================================================================
\section{Layer 2: The Spectral Action on $\mathbb{C}P^2 \times S^3$}
\label{sec:layer2}
%=============================================================================

%---------------------------------------------------------------------
\subsection{The Spectral Action Principle}
%---------------------------------------------------------------------

The DFD spectral action follows the Chamseddine--Connes framework:
\begin{equation}
S = \mathrm{Tr}\,f(P/\Lambda^2)\,,
\end{equation}
where $P$ is a Laplace-type operator on the internal space and $f$
is a cutoff function. The asymptotic expansion gives:
\begin{equation}
S \sim \sum_{n \ge 0} f_n \cdot a_n(P/\Lambda^2)\,,
\end{equation}
where $f_n$ are the moments of $f$ and $a_n$ are the Seeley--DeWitt
heat-kernel coefficients.

The gauge-kinetic terms arise from the $a_4$ coefficient, which
contains $F_{\mu\nu}^2$ terms. The coefficient of the $U(1)$ field
strength $B_{\mu\nu}^2$ determines $\kappa_{U(1)}$.

Source: \texttt{section\_alpha\_locked.tex}, Sec.~8C.2--8C.3.

%---------------------------------------------------------------------
\subsection{The Eight Locked Components (Table~XXVIII)}
%---------------------------------------------------------------------

The spectral action computation uses exactly eight inputs, all fixed
by geometry or Standard Model content:

\begin{table}[h]
\centering
\caption{Complete derivation chain for $\alpha^{-1}$
(Table~XXVIII of the Unified Theory, Sec.~8C).}
\label{tab:chain}
\renewcommand{\arraystretch}{1.3}
\begin{tabular}{llll}
\toprule
\textbf{Component} & \textbf{Value} & \textbf{Source} & \textbf{Status} \\
\midrule
$K_{\mathbb{C}P^2} = \mathcal{O}(-3)$ & $-3$
  & Algebraic geometry & Rigorous \\
$L_{\det} = K^{-1}$ & $\mathcal{O}(3)$
  & Spin$^c$ structure & Rigorous \\
$d = k_{\max} + 4$ & $64$
  & $\dim H^0(\mathcal{O}(k_{\max}+3))$ & Rigorous \\
$(d-1)/d$ & $63/64$
  & Traceless projection & Derived \\
$N_{\mathrm{species}}$ & $7$
  & SM SU(2) doublet components & SM content \\
$\mathrm{Tr}(Y^2)$ & $10$
  & SM hypercharge sum & SM content \\
$g_F$ & $8$
  & Spectral triple ($J \times \gamma \times \mathbb{C}$) & Derived \\
$w = N_{\mathrm{species}}/(g_F \cdot \mathrm{Tr}(Y^2))$ & $7/80$
  & Hypercharge weighting & Derived \\
$\varepsilon_{\mathrm{adj}} = d^2/(d^2-1)$ & $4096/4095$
  & Regular-module BOOST & Forced \\
$\beta_{U(1)}$ & $3.7969$
  & CS weighted average at $k_{\max}=60$ & Computed \\
\bottomrule
\end{tabular}
\end{table}

%---------------------------------------------------------------------
\subsection{Component 1: The Internal Space Geometry}
%---------------------------------------------------------------------

The internal manifold is $X = \mathbb{C}P^2 \times S^3$, a
7-dimensional compact space. Key geometric data:

\begin{itemize}
\item $\mathbb{C}P^2$ is a compact K\"ahler 4-manifold with
  $H^2(\mathbb{C}P^2;\mathbb{Z}) = \mathbb{Z}$
\item Canonical bundle: $K_{\mathbb{C}P^2} = \mathcal{O}(-3)$
  (the anti-canonical class is $3H$)
\item Spin$^c$ structure: $L_{\det} = K^{-1} = \mathcal{O}(3)$
  with $c_1(L_{\det}) = 3H$
\item $S^3$ supports SU(2) Chern--Simons theory; winding number
  in $\pi_3(S^3) = \mathbb{Z}$ gives baryon number
\end{itemize}

%---------------------------------------------------------------------
\subsection{Component 2: The Operator Choice (Locked)}
%---------------------------------------------------------------------

The operator is the \emph{connection Laplacian}:
\begin{equation}
P = -g^{ij}\nabla_i\nabla_j\,,
\end{equation}
acting on the internal bundle that carries the emergent gauge degrees
of freedom. The U(1) factor is implemented via twisting by a line
bundle over $\mathbb{C}P^2$ with curvature proportional to the
K\"ahler form $\omega$.

\paragraph{Why this is locked.}
The K\"ahler form is parallel ($\nabla\omega = 0$), so derivative
terms in higher Seeley--DeWitt coefficients vanish automatically.
The gauge-kinetic extraction from $a_4$ is unambiguous with this
operator choice.

Source: \texttt{section\_alpha\_locked.tex}, Sec.~8C.2.

%---------------------------------------------------------------------
\subsection{Component 3: The Plateau Cutoff (Locked)}
%---------------------------------------------------------------------

The cutoff function $f$ satisfies the \emph{plateau condition}:
$f^{(n)}(0) = 0$ for all $n \ge 1$, so all negative moments vanish:
\begin{equation}
f_{-2} = f_{-4} = \cdots = 0\,.
\end{equation}

\paragraph{Why this is locked.}
This eliminates $a_6$ and higher heat-kernel terms as hidden
ppm-level tuning knobs. With generic smooth cutoffs (e.g.\ Gaussian),
the $a_6$ contribution would be $\sim 2\%$---far too large and requiring
fine-tuned cancellation. The plateau cutoff is the unique choice that:
\begin{enumerate}
\item Preserves the leading $a_4$ gauge kinetic term
\item Eliminates subleading heat-kernel contributions
\item Requires no moment-tuning
\end{enumerate}

Source: \texttt{section\_alpha\_locked.tex}, Sec.~8C.3.

%---------------------------------------------------------------------
\subsection{Component 4: Toeplitz Truncation and the $(k+3)/(k+4)$ Factor}
%---------------------------------------------------------------------

The finite-$k$ truncation is implemented via Toeplitz quantization at
level $m = k + 3$ on $\mathbb{C}P^1 \subset \mathbb{C}P^2$:
\begin{equation}
d = \dim H^0(\mathbb{C}P^1, \mathcal{O}(m)) = m + 1 = k + 4\,.
\end{equation}

\paragraph{Origin of the $+3$ shift.}
The shift $m = k + 3$ arises from the Spin$^c$ structure:
\begin{equation}
K_{\mathbb{C}P^2} = \mathcal{O}(-3) \quad \Rightarrow \quad
L_{\det} = K^{-1} = \mathcal{O}(3)\,.
\end{equation}
When restricting to $\mathbb{C}P^1 \subset \mathbb{C}P^2$, the
line bundle $\mathcal{O}(k) \otimes L_{\det}$ becomes
$\mathcal{O}(k+3)$, giving sections of dimension $k + 4$.

\paragraph{The spectral cutoff.}
The determinant-channel removal at finite $d$ fixes:
\begin{equation}
\Lambda^3 = k \cdot \frac{d-1}{d} = k \cdot \frac{k+3}{k+4}\,.
\end{equation}

At $k_{\max} = 60$, $d = 64$:
\begin{equation}
\Lambda^3 = 60 \times \frac{63}{64} = 59.0625\,.
\end{equation}

The baseline normalization $\Lambda^3 = 885.9375$ (from $k = 60$,
$a = 9$, $n = 5$, $N_{\mathrm{gen}} = 3$) sets the overall scale.

Source: \texttt{section\_alpha\_locked.tex}, Sec.~8C.4.

%---------------------------------------------------------------------
\subsection{Component 5: The SM Content Inputs}
%---------------------------------------------------------------------

The spectral action $a_4$ trace over the internal fermion Hilbert
space requires three Standard Model inputs:

\begin{enumerate}
\item $N_{\mathrm{species}} = 7$: the number of SU(2) doublet
  components in one SM generation (counting $Q_L$ as 6 color-doublet
  components plus $L_L$ as 1 lepton doublet component = 7)

\item $\mathrm{Tr}(Y^2) = 10$: the hypercharge-squared trace over
  one SM generation:
  \[
  \mathrm{Tr}(Y^2) = \sum_R d_3 d_2 Y^2
  = 3 \cdot 2 \cdot (1/6)^2 + 3 \cdot 1 \cdot (2/3)^2
  + 3 \cdot 1 \cdot (1/3)^2 + 1 \cdot 2 \cdot (1/2)^2
  + 1 \cdot 1 \cdot 1^2 = 10/3 \times 3 = 10
  \]
  (using the one-generation sum with multiplicity $N_{\mathrm{gen}} = 3$
  included, or $10/3$ per generation)

\item $g_F = 8$: the spectral triple grading factor from real
  structure $J$, chirality $\gamma$, and complex structure
  $\mathbb{C}$ on the finite Hilbert space
\end{enumerate}

These combine into the hypercharge weighting:
\begin{equation}
w = \frac{N_{\mathrm{species}}}{g_F \cdot \mathrm{Tr}(Y^2)}
= \frac{7}{8 \times 10} = \frac{7}{80}\,.
\end{equation}

%---------------------------------------------------------------------
\subsection{Component 6: The Microsector Trace (BOOST Factor)}
\label{sec:boost}
%---------------------------------------------------------------------

This is the most consequential component. It involves a forced binary
fork between two possible trace normalizations.

\paragraph{Branch A: Regular-Module Microsector (survives).}
The finite Hilbert space is the algebra itself:
\begin{equation}
\mathcal{H}_F := A = M_d(\mathbb{C})\,,
\end{equation}
with democratic UV trace $\mathrm{tr}_{\mathrm{dem}} = (1/d^2)\,\mathrm{Tr}$.
The conversion to physics (per-generator) normalization is forced:
\begin{equation}
\varepsilon_{\mathrm{adj}}^{(A)} = \frac{d^2}{d^2 - 1}
= \frac{4096}{4095} = 1.000244\ldots
\end{equation}

\paragraph{Branch B: Fermion-Representation Microsector (falsified).}
If instead $\mathcal{H}_F \cong \mathbb{C}^d$:
\begin{equation}
\varepsilon_{\mathrm{adj}}^{(B)} = \frac{d^2 - 1}{d^2}
= \frac{4095}{4096} = 0.999756\ldots
\end{equation}

\paragraph{The fork resolves as follows:}
\begin{center}
\renewcommand{\arraystretch}{1.3}
\begin{tabular}{lccc}
\toprule
Branch & Factor & $\alpha^{-1}$ & Residual (ppm) \\
\midrule
\textbf{A (regular-module)} & $4096/4095$ & $137.03599985$ & $-0.006$ \\
B (fermion-rep) & $4095/4096$ & $137.03014445$ & $+42.7$ \\
\midrule
Experimental & --- & $137.035999084$ & --- \\
\bottomrule
\end{tabular}
\end{center}

Branch B misses by 43~ppm and cannot be rescued by adjusting any other
component. Under the no-hidden-knobs policy, only Branch A survives.

Source: \texttt{section\_alpha\_locked.tex}, Sec.~8C.5--8C.6.

%---------------------------------------------------------------------
\subsection{Component 7: The Chern--Simons Level Sum ($\beta_{U(1)}$)}
%---------------------------------------------------------------------

The Chern--Simons theory on $S^3$ provides the bare U(1) coupling via
a weighted sum over levels:
\begin{equation}
\beta_{U(1)} = \langle k + 2 \rangle
= \frac{\sum_{k=0}^{k_{\max}-1}(k+2)\,w(k)}
       {\sum_{k=0}^{k_{\max}-1} w(k)}\,,
\end{equation}
where the SU(2) Chern--Simons weights are:
\begin{equation}
w(k) = \frac{2}{k+2}\sin^2\frac{\pi}{k+2}\,.
\end{equation}

The quantum level shift $k \to k + h^\vee$ with $h^\vee(\mathrm{SU}(2)) = 2$
produces the $(k+2)$ factor (standard WZW/CS result).

\paragraph{Numerical evaluation at $k_{\max} = 60$:}
\begin{equation}
\beta_{U(1)}\big|_{k_{\max}=60}
= \frac{\sum_{k=0}^{59}(k+2)\,w(k)}{\sum_{k=0}^{59}w(k)}
= \frac{8.3245}{2.1924} = 3.796945\ldots \approx 3.80\,.
\end{equation}

The Wilson ratio $\beta_{SU(2)}/\beta_{U(1)} = (n_2/n_1) \times
N_{\mathrm{gen}} = 2 \times 3 = 6$ then gives
$\beta_{SU(2)} = 22.80$.

Source: \texttt{appendix\_K.tex}, Sec.~K.1.

%---------------------------------------------------------------------
\subsection{Component 8: Determination of $k_{\max} = 60$}
%---------------------------------------------------------------------

The UV cutoff $k_{\max}$ is determined by the closed Spin$^c$ index
on $\mathbb{C}P^2$:
\begin{equation}
k_{\max} = \chi(\mathbb{C}P^2, E)\,,
\end{equation}
where the twist bundle is:
\begin{equation}
E = \mathcal{O}(9) \oplus \mathcal{O}^{\oplus 5}\,.
\end{equation}

The two integers have physical origins:
\begin{itemize}
\item $a = 9$: the minimal globally well-defined hypercharge twist.
  With $q_1 = 3$ (uniquely minimal U(1) flux quantum from Spin$^c$
  integrality, Lemma~F.3), the triple tensor power gives
  $L_Y^{\otimes 3} = \mathcal{O}(3)^{\otimes 3} = \mathcal{O}(9)$.

\item $n = 5$: the five hypercharged chiral multiplet types per
  generation: $\{Q, u^c, d^c, L, e^c\}$.
\end{itemize}

The holomorphic Euler characteristic:
\begin{equation}
\chi(\mathcal{O}(9)) = \binom{11}{2} = 55\,,\qquad
\chi(\mathcal{O}) = 1\,,
\end{equation}
gives:
\begin{equation}
\boxed{k_{\max} = 55 + 5 = 60\,.}
\end{equation}

The derivation chain is:
\begin{equation}
\text{SM hypercharges} \to q_1 = 3 \to a = 9
\to k_{\max} = 60 \to \alpha^{-1} = 137.036\,.
\end{equation}
Crucially, $\alpha$ appears only at the \emph{end} as output, not input.

Source: \texttt{appendix\_F.tex}, Lemmas~F.3--F.4;
\texttt{appendix\_K.tex}, Sec.~K.1.

%=============================================================================
\section{The Explicit Computation of $\kappa_{U(1)}$}
\label{sec:computation}
%=============================================================================

%---------------------------------------------------------------------
\subsection{The $a_4$ Seeley--DeWitt Coefficient}
%---------------------------------------------------------------------

The gauge-kinetic action emerges from the $a_4$ heat-kernel coefficient
of the connection Laplacian $P = -g^{ij}\nabla_i\nabla_j$ on the
Toeplitz-truncated $\mathbb{C}P^2 \times S^3$.

With the plateau cutoff (eliminating $a_6$ and higher), only $a_4$
contributes to the gauge kinetic term:
\begin{equation}
S_{\mathrm{gauge}} = f_0 \cdot a_4(P/\Lambda^2)
\supset -\frac{\kappa_{U(1)}}{2}\,B_{\mu\nu}B^{\mu\nu}
  - \frac{\kappa_{SU(2)}}{2}\,W_{\mu\nu}^a W^{a\,\mu\nu}
  - \frac{\kappa_{SU(3)}}{2}\,G_{\mu\nu}^A G^{A\,\mu\nu}\,.
\end{equation}

The stiffness coefficients $\kappa_r$ absorb all the geometric
and SM trace factors from the internal space integration.

%---------------------------------------------------------------------
\subsection{Assembling $\kappa_{U(1)}$ from the Eight Components}
%---------------------------------------------------------------------

The computation of $\kappa_{U(1)}$ from the spectral action involves:

\begin{enumerate}
\item \textbf{Bare coupling from Chern--Simons:}
  $\beta_{U(1)} = 3.7969$ (from the level sum at $k_{\max} = 60$)

\item \textbf{Traceless projection:} The factor $(d-1)/d = 63/64$
  removes the determinant (trace) channel from $\mathrm{SU}(d)$

\item \textbf{SM spectral triple correction:} The hypercharge weighting
  $w = 7/80$ enters through the internal trace over the finite
  Hilbert space. Combined with $N_{\mathrm{gen}} = 3$ and $n_2 = 2$,
  this gives a correction factor
  $\mathcal{C}_{\mathrm{spectral}} = N_{\mathrm{species}}/(N_{\mathrm{gen}} \cdot n_2) = 7/6$

\item \textbf{BOOST factor:} The regular-module trace conversion
  $\varepsilon_{\mathrm{adj}} = 4096/4095$

\item \textbf{Electroweak mixing:} The factor $3/2$ from
  $1/e^2 = (3/2)\kappa_{U(1)}$ converts stiffness to $\alpha$
  via $\alpha^{-1} = 6\pi\,\kappa_{U(1)}$
\end{enumerate}

\paragraph{Key distinction.}
The computation of $\kappa_{U(1)}$ from the spectral action is
\emph{not} a single closed-form algebraic expression. It involves a
numerical evaluation of the $a_4$ Seeley--DeWitt coefficient with
all eight locked inputs. The closest approximation to a closed form
is:
\begin{equation}
\alpha^{-1} \approx \frac{35\pi}{3}\;\beta_{U(1)}\;\frac{d-1}{d}
\;\frac{d^2}{d^2-1}\,,
\end{equation}
which gives $137.024$ (85~ppm residual). The factors decompose as:
\begin{equation}
\frac{35\pi}{3} = 4\pi \times \underbrace{\frac{5}{2}}_{\text{EW mixing}}
\times \underbrace{\frac{7}{6}}_{\text{spectral triple}}\,.
\end{equation}

The full sub-ppm result requires the complete spectral action
evaluation.

Source: \texttt{Alpha\_Inverse\_Formula.tex}, Secs.~3--5.

%---------------------------------------------------------------------
\subsection{The Final Numerical Result}
%---------------------------------------------------------------------

\begin{equation}
\kappa_{U(1)}^{\mathrm{pred}} = 7.26999\ldots
\end{equation}

Therefore:
\begin{equation}
\alpha^{-1} = 6\pi \times 7.26999\ldots = 137.03599985\,.
\end{equation}

Residual from experiment:
\begin{equation}
\frac{\alpha^{-1}_{\mathrm{DFD}} - \alpha^{-1}_{\mathrm{exp}}}
     {\alpha^{-1}_{\mathrm{exp}}}
= \frac{137.03599985 - 137.035999084}{137.035999084}
= -0.006\text{~ppm}\,.
\end{equation}

%=============================================================================
\section{Lattice Monte Carlo Verification ($\sim 1\%$ Precision)}
\label{sec:lattice}
%=============================================================================

Independent of the spectral action route, lattice Monte Carlo
simulations verify the prediction at $\sim 1\%$ precision.

At the derived parameter point
$(\beta_{U(1)}, \beta_{SU(2)}) = (3.80, 22.80)$,
lattice simulations measure the renormalized stiffnesses via the
background-field method:

\begin{center}
\renewcommand{\arraystretch}{1.2}
\begin{tabular}{lccc}
\toprule
$L$ & $\alpha_W$ (mean) & $\sigma_\alpha$ & $\Delta\alpha/\alpha$ \\
\midrule
6 & 0.007297 & $9.4 \times 10^{-5}$ & $-0.00\%$ \\
8 & 0.007322 & $9.5 \times 10^{-5}$ & $+0.34\%$ \\
10 & 0.007361 & $6.8 \times 10^{-5}$ & $+0.88\%$ \\
12 & 0.007291 & $2.2 \times 10^{-5}$ & $-0.08\%$ \\
16 & 0.007380 & $1.1 \times 10^{-4}$ & $+1.13\%$ \\
\bottomrule
\end{tabular}
\end{center}

Key result (86 lattice runs, $L = 4$--$16$):
\begin{equation}
\alpha_W = 0.007297 \pm 0.000094
\quad\Longrightarrow\quad
\alpha^{-1} = 137.0 \pm 1.8
\quad(\sim 1\%\text{ precision}).
\end{equation}

From L12 data: $\kappa_{U(1)} \approx 7.27$,
$\kappa_{SU(2)} \approx 14.7$, giving
$\alpha^{-1} = 6\pi \times 7.27 = 137.1$.

Source: \texttt{appendix\_K.tex}, Sec.~K.2.

%=============================================================================
\section{The Additive Structure of the BOOST}
\label{sec:additive}
%=============================================================================

The two microsector branches decompose as:
\begin{equation}
\alpha^{-1} = F + G \cdot \delta\varepsilon\,,
\end{equation}
where $F \approx 137.033$, $G \approx 12$, and:
\begin{itemize}
\item Branch A: $\delta\varepsilon = +1/(d^2 - 1) = +1/4095$
\item Branch B: $\delta\varepsilon = -1/d^2 = -1/4096$
\end{itemize}

The coefficient $G \approx 12 = 4 \times N_{\mathrm{gen}}$ reflects
the three-generation structure of the microsector trace.

The BOOST factor $4096/4095$ thus contributes approximately
$12/4095 \approx 0.003$ to $\alpha^{-1}$, shifting the result from
$\sim 137.033$ to $137.036$---the difference between a 43~ppm miss
and a sub-ppm match.

%=============================================================================
\section{Relation to Standard Spectral Action Literature}
\label{sec:comparison}
%=============================================================================

%---------------------------------------------------------------------
\subsection{Comparison with Chamseddine--Connes}
%---------------------------------------------------------------------

In the standard Chamseddine--Connes spectral action on $M^4 \times F$
(with $F$ the finite spectral triple of the Standard Model), the
asymptotic expansion of $\mathrm{Tr}[f(D^2/\Lambda^2)]$ gives gauge
kinetic terms:
\begin{equation}
S_{\mathrm{gauge}}^{\mathrm{CC}} = \frac{f_0}{2\pi^2}
\left[a\,G_{\mu\nu}^2 + b\,W_{\mu\nu}^2 + c\,B_{\mu\nu}^2\right]\,,
\end{equation}
where the coefficients $a$, $b$, $c$ depend on the trace over the
internal Hilbert space $\mathcal{H}_F$.

In the standard Chamseddine--Connes model, the coefficients $a$, $b$, $c$
are determined by the representation content and satisfy specific ratios
that constrain gauge coupling unification.

\paragraph{DFD modifications.}
DFD differs from the standard approach in three key ways:

\begin{enumerate}
\item \textbf{Internal geometry:} DFD uses $\mathbb{C}P^2 \times S^3$
  (a smooth 7-manifold) rather than the purely algebraic finite space $F$
  of Chamseddine--Connes. This provides a geometric origin for the
  gauge group.

\item \textbf{Toeplitz truncation:} The finite-dimensionality is
  implemented via Toeplitz quantization at level $d = 64$, rather than
  being an intrinsic property of a finite algebra. This introduces the
  $(d-1)/d = 63/64$ traceless projection factor.

\item \textbf{Microsector trace:} The regular-module choice
  $\mathcal{H}_F = M_d(\mathbb{C})$ with BOOST factor $4096/4095$
  replaces the fermion-representation space used in standard
  noncommutative geometry.
\end{enumerate}

%---------------------------------------------------------------------
\subsection{The DFD Analogue of Coefficients $a$, $b$, $c$}
%---------------------------------------------------------------------

In DFD, the role of the Chamseddine--Connes coefficients $a$, $b$, $c$
is played by the stiffness coefficients $\kappa_r$:

\begin{center}
\renewcommand{\arraystretch}{1.2}
\begin{tabular}{lcc}
\toprule
Gauge sector & CC coefficient & DFD stiffness \\
\midrule
SU(3) & $a$ & $\kappa_{SU(3)} = 3\kappa_0$ \\
SU(2) & $b$ & $\kappa_{SU(2)} = 2\kappa_0$ \\
U(1)  & $c$ & $\kappa_{U(1)} = \kappa_0 = 7.26999\ldots$ \\
\bottomrule
\end{tabular}
\end{center}

The ratio $\kappa_{U(1)} : \kappa_{SU(2)} : \kappa_{SU(3)} = 1 : 2 : 3$
is fixed by the Ricci curvature of the corresponding projective spaces
(Frame Stiffness Theorem, Thm.~F.6).

The absolute value $\kappa_0 = \kappa_{U(1)} = 7.26999\ldots$ is
determined by the spectral action computation with all eight locked
inputs.

%=============================================================================
\section{The Coefficient of $B_{\mu\nu}^2$: Explicit Identification}
\label{sec:Bmunu}
%=============================================================================

The question posed is: what is the explicit coefficient of $B_{\mu\nu}^2$
(the U(1) hypercharge field strength) in the DFD spectral action?

\begin{tcolorbox}[colback=green!5!white, colframe=green!50!black,
  title=The U(1) Stiffness Coefficient]
In the DFD spectral action on the Toeplitz-truncated
$\mathbb{C}P^2 \times S^3$ at $d = 64$:

\begin{equation}
\boxed{\kappa_{U(1)} = 7.26999\ldots = \frac{\alpha^{-1}_{\mathrm{exp}}}{6\pi}}
\end{equation}

This is the coefficient of $-\frac{1}{2}B_{\mu\nu}B^{\mu\nu}$ in the
gauge-kinetic action extracted from the $a_4$ Seeley--DeWitt coefficient.

The other stiffnesses follow from the Frame Stiffness Theorem:
\begin{align}
\kappa_{SU(2)} &= 2\kappa_{U(1)} = 14.5400\ldots \\
\kappa_{SU(3)} &= 3\kappa_{U(1)} = 21.8100\ldots
\end{align}
\end{tcolorbox}

\paragraph{Why $\kappa_{U(1)} = 7.26999$ specifically.}
The value arises from the interplay of:
\begin{itemize}
\item The Chern--Simons bare coupling $\beta_{U(1)} = 3.7969$
  (from the $k_{\max} = 60$ level sum)
\item The traceless projection $(d-1)/d = 63/64$
\item The SM spectral triple correction factor $\sim 7/6$
\item The BOOST factor $4096/4095$
\item The electroweak mixing factor $3/2$
  (absorbed into the $6\pi$ of the master formula)
\end{itemize}

The approximate relation is:
\begin{equation}
\kappa_{U(1)} \approx \frac{35}{18}\,\beta_{U(1)}\,\frac{63}{64}
\,\frac{4096}{4095}
= \frac{35}{18} \times 3.7969 \times 0.984375 \times 1.000244
= 7.268\ldots
\end{equation}
(within 85~ppm of the exact value; the full spectral action evaluation
is needed for sub-ppm precision).

%=============================================================================
\section{The Complete Derivation Chain (Diagram)}
\label{sec:diagram}
%=============================================================================

\begin{center}
\fbox{\parbox{0.92\textwidth}{
\textbf{SM hypercharge spectrum} $\to$ $q_1 = 3$ (anomaly cancellation)

$\downarrow$

$a = q_1^2 \cdot q_1 = 9$ $\to$ twist bundle $E = \mathcal{O}(9) \oplus \mathcal{O}^{\oplus 5}$

$\downarrow$

$k_{\max} = \chi(\mathbb{C}P^2, E) = 55 + 5 = 60$ \quad [Bridge Lemma]

$\downarrow$

$\beta_{U(1)} = \langle k+2 \rangle_{k_{\max}=60} = 3.7969$ \quad [CS weighted average]

$\downarrow$ \quad + \quad Wilson ratio $= 6$ \quad +
  \quad $\kappa_{U(1)}/\kappa_{SU(2)} = 1/2$

$\downarrow$

\textbf{Spectral action route:} $a_4$ on Toeplitz-truncated
$\mathbb{C}P^2 \times S^3$

+ traceless $(63/64)$ + BOOST $(4096/4095)$
+ SM content $(N_{\mathrm{species}}, g_F, \mathrm{Tr}(Y^2))$

$\to$ $\kappa_{U(1)} = 7.26999\ldots$ $\to$
$\alpha^{-1} = 6\pi\kappa = 137.03599985$ \quad (sub-ppm)

\medskip
\textbf{OR (independent verification):}

\textbf{Lattice route:} MC at $(\beta_{U(1)}, \beta_{SU(2)}) = (3.80, 22.80)$

$\to$ measure $\kappa_{U(1)} \approx 7.27$ $\to$
$\alpha^{-1} = 6\pi\kappa \approx 137$ \quad ($\pm 1\%$)
}}
\end{center}

%=============================================================================
\section{Summary and Key Numbers}
\label{sec:summary}
%=============================================================================

\begin{tcolorbox}[colback=blue!5!white, colframe=blue!50!black,
  title=Summary: The Spectral Action Computation in DFD]

\textbf{Master formula (exact):}
\begin{equation}
\alpha^{-1} = 6\pi\,\kappa_{U(1)}
\end{equation}

\textbf{Key numerical values:}
\begin{center}
\renewcommand{\arraystretch}{1.2}
\begin{tabular}{ll}
\toprule
Quantity & Value \\
\midrule
$k_{\max}$ (UV cutoff) & $60$ \\
$d$ (Toeplitz dimension) & $64$ \\
$(d-1)/d$ (traceless projection) & $63/64 = 0.984375$ \\
$\varepsilon_{\mathrm{adj}}$ (BOOST) & $4096/4095 = 1.000244\ldots$ \\
$\beta_{U(1)}$ (CS bare coupling) & $3.7969$ \\
$w$ (hypercharge weight) & $7/80 = 0.0875$ \\
$\kappa_{U(1)}$ (U(1) stiffness) & $7.26999\ldots$ \\
$\kappa_{SU(2)}$ (SU(2) stiffness) & $14.5400\ldots$ \\
$\alpha^{-1}$ (DFD prediction) & $137.03599985$ \\
$\alpha^{-1}$ (experiment) & $137.035999084(21)$ \\
Residual & $-0.006$~ppm \\
\bottomrule
\end{tabular}
\end{center}

\textbf{Locked conventions (no free parameters):}
\begin{itemize}
\item Operator: connection Laplacian with parallel curvature
\item Regulator: plateau cutoff ($f_{-2} = f_{-4} = \cdots = 0$)
\item Finite-$k$: Toeplitz truncation with $d = k + 4 = 64$
\item Microsector: regular-module ($\mathcal{H}_F = M_d(\mathbb{C})$)
\item Trace: democratic UV $\to$ per-generator physics (BOOST forced)
\end{itemize}

\textbf{The fermion-rep microsector is falsified:}
43~ppm deficit, unrescuable under no-knobs policy.
\end{tcolorbox}

%=============================================================================
\section*{Source File Index}
%=============================================================================

All referenced files are in the DFD codebase at:\\
\texttt{/Users/garyalcock/claudecode/densityfielddynamics/}

\begin{center}
\small
\renewcommand{\arraystretch}{1.1}
\begin{tabular}{ll}
\toprule
\textbf{File} & \textbf{Content} \\
\midrule
\texttt{..\_108.../section\_alpha\_locked.tex}
  & Convention-locked $\alpha$ (Sec.~8C) \\
\texttt{..\_108.../section\_alpha\_relations.tex}
  & The $\alpha$-relations (Sec.~8) \\
\texttt{..\_108.../section\_gauge\_coupling.tex}
  & Gauge coupling variation (Sec.~8B) \\
\texttt{..\_108.../appendix\_F.tex}
  & Gauge emergence proofs \\
\texttt{..\_108.../appendix\_G.tex}
  & $\alpha$-relation derivations \\
\texttt{..\_108.../appendix\_K.tex}
  & Microsector physics, CS derivation, lattice \\
\texttt{..\_108.../appendix\_O\_alpha57.tex}
  & $\alpha^{57}$ exponent \\
\texttt{..\_108.../supplementary/PRL\_SM\_from\_Topology.tex}
  & PRL-style SM from topology \\
\texttt{DFD\_Research\_Output/.../Alpha\_Inverse\_Formula.tex}
  & Explicit chain document \\
\bottomrule
\end{tabular}
\end{center}

\end{document}
